\emph{(i)} Theorem \claimref{synthesis.de-sitter-entropy-pi-n-squared}.

\emph{(ii)} Proposition \claimref{scrambling.residue-each-convention} writes the count registered per modular period as $N_{\mathrm{epoch}} = \kappa_{c} \times A \times 1$, horizon-normalised cells times one tick per modular period. Under the hypothesis that the cosmological horizon's modes are counted by the same rule, $|J| = \kappa_{c}A_{\Lambda} = 4\pi\kappa_{c}N^{2}$, which is $N^{2}$ times a constant for every $\kappa_{c}$. The residues are those of the same Proposition: $\kappa_{c} = 1/(4\ln 2)$ in the capacity convention, its clause (ii), gives $|J| = \pi N^{2}/\ln 2$ and $|J|\ln 2 = \pi N^{2} = S_{\Lambda}$; $\kappa_{c} = \tfrac14$ in the 1-nat convention, its clause (i), gives $|J| = \pi N^{2} = S_{\Lambda}$. The rule counts per modular period, and the cosmological horizon's modular period is $2\pi/H$ under (P) and (R$_{\Lambda}$); the count itself does not use the period's value.

\emph{(iii)} From (ii) in the capacity convention, $C = |J|\ln 2 = \pi N^{2}$, so $N = \sqrt{C/\pi}$.
